\documentclass[11pt,a4paper]{article}

\usepackage[T1]{fontenc}
\usepackage[utf8]{inputenc}
\usepackage{geometry}
\usepackage{array}
\usepackage{booktabs}
\usepackage{longtable}
\usepackage{xcolor}
\usepackage{hyperref}
\usepackage{listings}

\geometry{margin=2.4cm}
\hypersetup{
  colorlinks=true,
  linkcolor=blue!50!black,
  urlcolor=blue!60!black,
  citecolor=blue!50!black
}
\lstset{
  basicstyle=\ttfamily\small,
  columns=fullflexible,
  breaklines=true,
  frame=single,
  framerule=0.3pt,
  backgroundcolor=\color{gray!5}
}

\title{ILB Incubator Management Platform\\\large Final Assignment Report}
\author{SDL Project Group 20}
\date{May 2026}

\begin{document}
\maketitle

\begin{abstract}
This report summarises the final delivery of the ILB Incubator Management
Platform: a Docker Compose deployable microservices system for managing
incubated companies, contracts, finance, spaces, bookings, inventory, support
tickets, and documents. The final branch consolidates strict service boundaries,
role-aware frontend workflows, RabbitMQ event integration, operational
runbooks, and a GitLab CI/CD path for variable-driven production deployment.
\end{abstract}

\tableofcontents
\newpage

\section{Project scope}
The platform supports incubator staff, client-company users, and public booking
requesters. Staff users operate the internal dashboard, companies, contracts,
finance, spaces, bookings, inventory, tickets, users, and documents. Client
users access company-scoped contracts, payments, bookings, tickets, and company
profile information. Public users can submit booking requests without an
authenticated session.

The final objective was to deliver a production-shaped system that can be run
locally with Docker Compose, explained during the defense, and promoted through
GitLab CI/CD without committing environment-specific secrets.

\section{Architecture overview}
The repository is a monorepo, but the runtime model is deliberately
microservice-oriented. Each bounded context is a separate Django application and
owns its own PostgreSQL database. Services communicate through REST for
request/response flows and RabbitMQ events for asynchronous consistency.

\begin{longtable}{p{0.24\textwidth}p{0.36\textwidth}p{0.22\textwidth}p{0.12\textwidth}}
\toprule
\textbf{Compose service} & \textbf{Bounded context} & \textbf{Gateway prefix} & \textbf{Port} \\
\midrule
\endhead
auth-service & Identity, sessions, JWT issuance and validation & /api/auth/ & 8001 \\
company-service & Incubated company profiles, CAE, maturity and archive state & /api/companies/ & 8002 \\
contract-service & Company contracts and lifecycle & /api/contracts/ & 8003 \\
finance-service & Payments, billing, overdue processing, reports & /api/finance/ & 8004 \\
space-service & Physical spaces, capacity and occupancy projections & /api/spaces/ & 8005 \\
booking-service & Staff, client and public booking workflows & /api/bookings/ & 8006 \\
inventory-service & Equipment catalogue, assignment and release & /api/inventory/ & 8007 \\
ticket-service & Support queues and ticket threads & /api/tickets/ & 8008 \\
dashboard-service & Cross-service operational aggregates & /api/dashboard/ & 8009 \\
document-service & Document metadata and MinIO-backed object storage & /api/documents/ & 8010 \\
\bottomrule
\end{longtable}

The Nginx gateway is the external HTTP boundary. It serves the Next.js frontend
and routes API traffic by path prefix. Authentication uses gateway validation
against the auth service and forwards trusted identity headers to downstream
services. Internal Django service ports remain private to the Compose network.

\section{Implemented capabilities}
\subsection{Backend services}
The final implementation covers the main demo path across all service domains:
\begin{itemize}
  \item Auth: login, JWT handling, user management, gateway verification and RBAC.
  \item Company: company CRUD, maturity/status data, employees, statistics and events.
  \item Contract: lifecycle endpoints, monthly fee snapshots, expiration handling and events.
  \item Finance: payments, billing/overdue commands, event consumers and reports.
  \item Space: type and space CRUD, occupancy, contract and booking projections.
  \item Booking: staff/client/public booking flows and lifecycle events.
  \item Inventory: equipment/type CRUD, assignment/release and booking event handlers.
  \item Ticket: scoped tickets, message threads and staff metrics.
  \item Dashboard: service health, ticket KPIs and aggregate snapshots.
  \item Document: metadata plus MinIO-backed upload, download, list and delete flows.
\end{itemize}

\subsection{Frontend workflows}
The frontend is a Next.js 14 App Router application with TypeScript and Ant
Design. It separates staff routes from client portal routes and keeps the public
booking request outside the authenticated shell. The delivered pages support the
defense walkthrough: staff dashboard, company/contract/finance/spaces/bookings/
inventory/tickets/users areas, client portal pages, and public booking request.

\subsection{Event-driven integration}
Cross-service consistency is implemented with the \texttt{incubator.events}
RabbitMQ topic exchange and a shared envelope from \texttt{libs/py-common}.
Events include \texttt{contract.activated}, \texttt{contract.terminated},
\texttt{booking.approved}, \texttt{booking.cancelled},
\texttt{booking.completed}, and \texttt{payment.recorded}. Consumers persist or
check \texttt{event\_id} values so duplicate delivery can be retried safely.
Recurring work is implemented with scheduler sidecar containers that run Django
management commands from cron files; the design intentionally avoids Celery.

\section{Runtime and deployment}
\subsection{Docker Compose runtime}
The production-shaped runtime is defined in \texttt{infra/docker-compose.yml}.
It includes PostgreSQL, Redis, RabbitMQ, MinIO, one Django container per bounded
context, scheduler/consumer sidecars where needed, the Next.js frontend, and the
Nginx gateway. Environment variables are sourced from a root \texttt{.env} file;
\texttt{.env.example} documents the required non-secret defaults and values to
review before deployment.

\subsection{GitLab CI/CD production deployment}
The final branch includes a variable-driven production deployment path in
\texttt{.gitlab-ci.yml} and \texttt{infra/scripts/ci/deploy-remote.sh}. The
manual \texttt{deploy:production} job runs on \texttt{main}, uses protected
GitLab variables for SSH access and the production environment file, archives
the current commit to the server, validates the Compose configuration, and runs
Docker Compose remotely.

Changed-service selection is handled by
\texttt{infra/scripts/ci/detect-changed-services.sh}. Service-local changes rebuild only
the affected Compose services and sidecars. Shared library, Compose, deploy
script, or environment-template changes trigger a full application rollout to
avoid stale dependent images.

\begin{lstlisting}[language=bash,caption={Production deploy command shape}]
docker compose --env-file .env --project-directory . \
  -f infra/docker-compose.yml config --quiet

docker compose --env-file .env --project-directory . \
  -f infra/docker-compose.yml up -d --build --remove-orphans \
  <changed-services>
\end{lstlisting}

Required GitLab variables are \texttt{PRODUCTION\_DEPLOY\_HOST},
\texttt{PRODUCTION\_DEPLOY\_USER}, \texttt{PRODUCTION\_DEPLOY\_PATH},
\texttt{PRODUCTION\_DEPLOY\_SSH\_PRIVATE\_KEY}, and either
\texttt{PRODUCTION\_ENV\_FILE\_BASE64} or \texttt{PRODUCTION\_ENV\_FILE}.
No production hostnames, keys, or secrets are committed to the repository.

\section{Quality evidence}
The repository contains defense and release evidence under
\texttt{docs/defense/}. The latest recorded release criteria require host-only
quality gates, Docker-backed Compose evidence on a host with Docker socket
access, frontend typecheck/unit/lint/build, browser smoke, and updated operator
documentation.

For this final delivery task, the following checks were rerun in the worker
branch:
\begin{itemize}
  \item \texttt{bash -n infra/scripts/ci/*.sh}
  \item \texttt{docker-compose -f infra/docker-compose.yml --project-directory . config --quiet}
  \item GitLab YAML parse for \texttt{.gitlab-ci.yml} and child CI files.
  \item Changed-service selector smoke tests in a temporary Git repository.
  \item \texttt{npx prettier --check .gitlab-ci.yml docs/deploy.md}
  \item \texttt{npm --prefix frontend run typecheck}
  \item \texttt{python3 -m pytest infra/tests -q} with 15 passing tests.
  \item \texttt{latexmk -pdf -interaction=nonstopmode -halt-on-error docs/final-report/final-report.tex}
\end{itemize}

\section{Known limitations and operational notes}
\begin{itemize}
  \item Live Compose and gateway Playwright evidence require a host that can
  access the Docker socket. If unavailable, the blocker must be recorded as an
  environment limitation rather than product evidence.
  \item The production deploy job was validated structurally and by script
  checks, but a live SSH deployment requires protected GitLab variables and a
  prepared Docker host.
  \item Dependency audit warnings are present in the existing frontend lockfile
  and should be handled in a separate dependency/security maintenance lane.
  \item Production secrets must stay in GitLab protected variables or a secrets
  manager; the repository should never contain real \texttt{.env} values.
\end{itemize}

\section{Conclusion}
The final assignment delivery provides a coherent, demonstrable incubator
management platform with strict service ownership, event-driven integration,
role-aware frontend workflows, Compose-based runtime operations, and a
variable-driven GitLab production deployment route. The included documentation,
defense scripts, CI/CD configuration, and compiled report support both technical
review and operational handoff.

\end{document}
